% Synthetic T2A regular/bold 10pt regression; no Unicode mapping patch.
\documentclass{article}
\usepackage[utf8]{inputenc}
\usepackage[T2A]{fontenc}
\usepackage[english,ukrainian]{babel}
\makeatletter
\typeout{ENGLISH-ID:\number\l@english}
\typeout{UKRAINIAN-ID:\number\l@ukrainian}
\makeatother
\begin{document}
\selectlanguage{ukrainian}
\typeout{UKRAINIAN-ACTIVE:\the\language}
\typeout{UKRAINIAN-CONTENTS:\meaning\contentsname}
\typeout{UKRAINIAN-CAPTION-EXPANSION:\meaning\ukrainiancontentsname}
\typeout{UKRAINIAN-HYPHENS-BEGIN}
\showhyphens{український українська університет комп'ютерний програмування}
\typeout{UKRAINIAN-HYPHENS-END}
\contentsname: український університет. Іі Її Єє Ґґ.
\textbf{Український текст. Іі Її Єє Ґґ.}
\selectlanguage{english}
\typeout{ENGLISH-ACTIVE:\the\language}
\typeout{ENGLISH-CONTENTS:\contentsname}
\typeout{ENGLISH-HYPHENS-BEGIN}
\showhyphens{hyphenation university compatibility}
\typeout{ENGLISH-HYPHENS-END}
\typeout{UKRAINIAN-IN-ENGLISH-CONTROL-BEGIN}
\showhyphens{\fontencoding{T2A}\selectfont український університет}
\typeout{UKRAINIAN-IN-ENGLISH-CONTROL-END}
ASCII I i 123. English retained.
\end{document}
